\documentclass[11pt,letterpaper]{article}

% ==============================================================================
%% Encoding & idioma
% ==============================================================================
\usepackage[utf8]{inputenc}
\usepackage[T1]{fontenc}
\usepackage[spanish]{babel}

% ==============================================================================
%% Tipografía y diseño
% ==============================================================================
\usepackage{lmodern}
\usepackage[margin=2.5cm]{geometry}
\usepackage{xcolor}
\usepackage{enumitem}
\usepackage{titlesec}
\usepackage{parskip}
\usepackage{array}
\usepackage{booktabs}
\usepackage{hyperref}

% ==============================================================================
%% Paleta de colores (consistente con plantilla Beamer Colmex)
% ==============================================================================
\definecolor{main}{HTML}{5E002B}
\definecolor{subtitle}{RGB}{75,75,75}
\definecolor{cerulean}{RGB}{0,123,167}
\definecolor{redorange}{RGB}{210,77,49}
\definecolor{forestgreen}{RGB}{36,128,103}

\hypersetup{
    colorlinks=true,
    linkcolor=main,
    citecolor=main,
    urlcolor=cerulean,
    pdftitle={Programación para Proyectos de Datos -- Syllabus},
    pdfauthor={Daniel Kelly}
}

% ==============================================================================
%% Estilo de secciones
% ==============================================================================
\titleformat{\section}
  {\Large\bfseries\color{main}}{\thesection}{1em}{}
\titleformat{\subsection}
  {\large\bfseries\color{main}}{\thesubsection}{1em}{}
\titleformat{\subsubsection}
  {\normalsize\bfseries\color{subtitle}}{}{0em}{}

% ==============================================================================
%% Comandos auxiliares (consistentes con plantilla Beamer)
% ==============================================================================
\newcommand{\blue}[1]{{\color{cerulean}#1}}
\newcommand{\orange}[1]{{\color{redorange}#1}}
\newcommand{\green}[1]{{\color{forestgreen}#1}}

% Marcador de contenido pendiente (para iterar sesión por sesión)
\newcommand{\porDesarrollar}{\textit{\color{redorange}[Por desarrollar]}}

% Bloque estandarizado para cada semana (cada semana contiene 2 sesiones de 1:30 hr)
\newcommand{\semana}[2]{%
  \subsubsection*{\color{main}Semana #1: #2}%
}

% Sub-bloque para cada sesión dentro de una semana
\newcommand{\subsesion}[2]{%
  \medskip\noindent\textbf{\color{main}Sesión #1 --- #2}\par\smallskip%
}

% ==============================================================================
%% Bibliografía
% ==============================================================================
\usepackage[style=apa, backend=biber]{biblatex}
\addbibresource{bibliography.bib}

% ==============================================================================
%% Encabezado
% ==============================================================================
\title{\color{main}\textbf{Programación para Proyectos de Datos}\\[0.3em]
       \large\color{subtitle}Syllabus}
\author{Daniel Kelly\\
        \small El Colegio de México --- Licenciatura en Economía}
\date{Otoño 2026}

\begin{document}
\maketitle

% ==============================================================================
\section{Presentación}
% ==============================================================================

Este curso complementa la formación cuantitativa de la Licenciatura en Economía con una visión integral del flujo de un proyecto de datos. Los cursos previos del programa equipan al estudiante con \blue{habilidades específicas} de programación y de análisis estadístico; el énfasis aquí está en cómo esas habilidades se \orange{articulan dentro de un mismo proyecto}: desde estructurar un repositorio reproducible y traer datos de fuentes diversas, hasta limpiarlos, modelarlos y producir un entregable comunicable.

La filosofía del curso es \textbf{\textit{zero-to-hero}}: se asume que el estudiante puede llegar con experiencia previa de programación limitada o nula, y termina capaz de levantar por su cuenta un proyecto de datos completo bajo prácticas reproducibles.

El lenguaje del curso es \textbf{R}. La elección se discute en la primera sesión; en términos generales, R es la lingua franca de la investigación cuantitativa académica y ofrece un ecosistema integrado para todas las etapas de un proyecto de datos. El IDE recomendado es \textbf{Positron}.

% ==============================================================================
\section{Objetivos}
% ==============================================================================

Al finalizar el curso, el estudiante será capaz de:

\begin{enumerate}[leftmargin=*]
    \item Estructurar un proyecto de datos completo bajo convenciones reproducibles (control de versiones, organización de directorios, gestión de dependencias).
    \item Importar, limpiar y transformar datos provenientes de fuentes heterogéneas (archivos planos, hojas de cálculo, bases de datos relacionales).
    \item Programar funciones reutilizables y aplicar paradigmas de programación funcional para automatizar tareas repetitivas.
    \item Estimar y reportar modelos estadísticos comunes en investigación económica desde el punto de vista programático.
    \item Producir entregables comunicables: visualizaciones, documentos automatizados, aplicaciones interactivas e integraciones con servicios externos.
\end{enumerate}

% ==============================================================================
\section{Prerrequisitos}
% ==============================================================================

\begin{itemize}[leftmargin=*]
    \item No se requiere experiencia previa en programación.
    \item Se asume manejo básico de estadística descriptiva y nociones iniciales de econometría (regresión lineal). El curso \textit{no} enseña la teoría de los modelos; solo cómo implementarlos.
\end{itemize}

% ==============================================================================
\section{Metodología}
% ==============================================================================

\textbf{Duración:} 16 semanas, con 2 sesiones de 1:30 hr por semana (32 sesiones totales).

\textbf{Estructura semanal.} Cada semana tiene dos sesiones con roles diferenciados:

\begin{itemize}[leftmargin=*]
    \item \textbf{\blue{Sesión 1 --- Teoría.}} Exposición conceptual del tema de la semana, con código en vivo a manera de ilustración. La proporción típica es de \textbf{$\approx$60\% teoría / 40\% código demostrativo}, aunque puede inclinarse hacia la teoría en semanas más conceptualmente densas.
    \item \textbf{\orange{Sesión 2 --- Práctica.}} Inicia con un \textit{resumen rápido} de los puntos clave de la sesión 1, y se dedica el tiempo restante a \textbf{ejercicios prácticos} resueltos en clase.
\end{itemize}

\textbf{Tareas.} El curso \textit{no} contempla tareas semanales fuera de clase. El trabajo evaluable se concentra en los ejercicios resueltos en la sesión 2 y en el \orange{proyecto integrador} que los estudiantes construyen a lo largo del semestre (ver sección correspondiente).

% ==============================================================================
\section{Evaluación}
% ==============================================================================

\porDesarrollar{} \emph{(porcentajes propuestos, por confirmar:)}

\begin{center}
\begin{tabular}{lr}
\toprule
\textbf{Componente} & \textbf{Peso} \\
\midrule
Checkpoints del proyecto integrador       & 40\% \\
Proyecto final                            & 30\% \\
Ejercicios en clase (sesión 2 semanal)    & 25\% \\
Asistencia                                & 5\%  \\
\bottomrule
\end{tabular}
\end{center}

% ==============================================================================
\section{Estructura del curso}
% ==============================================================================

El curso se organiza en cinco partes:

\begin{description}[leftmargin=2em, style=nextline]
    \item[\textcolor{main}{Parte I --- Fundamentos (semanas 1--4).}] Introducción a R, estructuras de datos, importación, exploración inicial y estructura reproducible de proyectos.
    \item[\textcolor{main}{Parte II --- Manipulación (semanas 5--7).}] Manejo de datos con \texttt{dplyr} y manipulación de texto.
    \item[\textcolor{main}{Parte III --- Programación (semanas 8--10).}] Control de flujo, funciones y programación funcional.
    \item[\textcolor{main}{Parte IV --- Análisis (semanas 11--12).}] Modelado estadístico (incluyendo series de tiempo a nivel programático) y datos a escala con SQL.
    \item[\textcolor{main}{Parte V --- Productos (semanas 13--16).}] Visualización con \texttt{ggplot2}, automatización de documentos, aplicaciones \texttt{Shiny} e interfaces de modelos de lenguaje.
\end{description}

% ==============================================================================
\section{Temario detallado}
% ==============================================================================

% ------------------------------------------------------------------------------
\subsection*{\color{main}Parte I --- Fundamentos}
% ------------------------------------------------------------------------------

\semana{1}{¿Por qué R? Setup y primeros pasos}

Esta semana es la entrada al curso. Cubre la motivación detrás de elegir R, la instalación y configuración del entorno de trabajo, los primeros conceptos del lenguaje (asignación, vectores, tipos), y la presentación formal del proyecto integrador. Es la única semana con peso significativo en logística administrativa.

\subsesion{1.1}{Teoría}

\begin{itemize}[leftmargin=*, itemsep=2pt]
    \item \textbf{Bienvenida y proyecto integrador} (15 min). Estructura del curso, esquema de evaluación, presentación formal del proyecto integrador y del pipeline ETL ($Extract$ $\to$ $Transform$ $\to$ $Load$ $\to$ análisis $\to$ producto).
    \item \textbf{¿Por qué R?} (15 min). R en el contexto de la investigación económica académica; contraste con Python; razones para descartar Stata. Ecosistema CRAN, tidyverse y Posit.
    \item \textbf{Uso responsable de modelos de lenguaje en el curso} (10 min). Diferencia operativa entre usar un LLM como \blue{chatbot} (pegar pregunta, copiar respuesta) y como \orange{agente integrado en el proyecto} (lee archivos, propone cambios, ejecuta comandos con el humano en el \textit{loop} ---por ejemplo, Claude Code o el modo agente de Cursor). Depender mal de un LLM \textbf{atrofia la capacidad individual} de razonar sobre código; usado como agente puede acelerar el trabajo de quien ya domina los fundamentos. \textbf{Política del curso:} agentes permitidos en el proyecto integrador; \textit{no} permitidos en los ejercicios de clase. El estudiante debe ser capaz de leer y modificar cualquier línea que entregue.
    \item \textbf{Características del lenguaje} (15 min). Interpretado, vectorizado, multiparadigma. Cómo se ejecuta R: REPL, script con \texttt{source()}, ejecutable con \texttt{Rscript}.
    \item \textbf{Setup de Positron} (15 min). Instalación, layout del IDE (consola, source, environment, plots, packages), configuración mínima.
    \item \textbf{Primeros objetos en R} (20 min). Asignación (\texttt{<-} vs \texttt{=}); tipos atómicos (\texttt{numeric}, \texttt{integer}, \texttt{logical}, \texttt{character}); vectores con \texttt{c()}, vectorización, operaciones aritméticas; funciones de inspección (\texttt{class}, \texttt{typeof}, \texttt{length}, \texttt{str}).
\end{itemize}

\subsesion{1.2}{Práctica}

\textbf{Resumen rápido} (15 min): hoja-resumen con asignación, vectores, tipos, funciones de inspección y operaciones aritméticas.

\textbf{Ejercicios en clase} (1:15 hr):
\begin{enumerate}[leftmargin=*, itemsep=2pt]
    \item \textbf{Verificación de setup.} Cada estudiante abre Positron, crea un proyecto y ejecuta \texttt{hola\_mundo.R}.
    \item \textbf{Calculadora vectorizada.} Dados vectores de precios y cantidades, calcular ingreso total, promedio y máximo. Introduce \texttt{sum()}, \texttt{mean()}, \texttt{max()}.
    \item \textbf{Coerción de tipos.} Construir \texttt{c(1, "a")}, \texttt{c(1, TRUE)}, \texttt{c(1L, 2.5)}; observar la jerarquía de coerción y discutir.
    \item \textbf{Mini-script con \texttt{Rscript}.} Crear un script con comentarios, asignaciones y \texttt{print()} final; correrlo desde la terminal con \texttt{Rscript}.
    \item \textbf{Cierre + proyecto integrador.} Cada estudiante declara su encuesta INEGI elegida, crea su repositorio Git con estructura \texttt{input/output/docs/files} y hace el primer \textit{commit}.
\end{enumerate}

\medskip
\noindent\textbf{Objetivos de aprendizaje}
\begin{enumerate}[leftmargin=*, itemsep=1pt]
    \item Justificar la elección de R sobre alternativas en el contexto de un proyecto académico de investigación económica.
    \item Tener Positron + R instalados y configurados, capaces de ejecutar un \textit{script} básico desde el IDE y desde la terminal.
    \item Crear, asignar y operar con vectores atómicos en R, identificando su tipo y tamaño.
    \item Distinguir el uso de un LLM como chatbot frente a su uso como agente integrado al proyecto, y la política del curso al respecto.
\end{enumerate}

\medskip
\noindent\textbf{Lecturas}
\begin{itemize}[leftmargin=*, itemsep=1pt]
    \item \textit{R for Data Science} (2.\textsuperscript{a} ed.), Cap.\ 1 (\textit{Introduction}) y Cap.\ 2 (\textit{Workflow: basics}) \cite{wickham2023r4ds}.
    \item \textit{Advanced R} (2.\textsuperscript{a} ed.), Cap.\ 2 §2.1--2.2 y Cap.\ 3 §3.1--3.2 \cite{wickham2019advr}.
\end{itemize}

\medskip
\noindent\textbf{Conexión con proyecto integrador}

Anuncio formal del proyecto en Sesión 1.1 y formalización de la elección de encuesta en Sesión 1.2. El proyecto es \textbf{individual}; cada estudiante elige una encuesta del INEGI sujeta a tres restricciones: (i) debe tener \textbf{microdatos}; (ii) debe contar con \textbf{al menos dos módulos relacionables} (e.g.\ ENSU con CS/CB/VIV; ENIGH con viviendas/hogares/personas/concentrado); y (iii) debe tener \textbf{al menos dos levantamientos en el tiempo}. El producto final será un \blue{Tablero Shiny} acompañado de un \blue{reporte periódico automatizado} con \texttt{officer}. \textbf{Entregable del Checkpoint 1} (cierre de la Parte I, Semana 4): repositorio Git con estructura \texttt{input/output/docs/files}, encuesta justificada en una página, descriptor técnico leído, y primer \textit{script} de extracción.

\semana{2}{Estructuras de datos e indexación}
\porDesarrollar

\semana{3}{Importación de datos y EDA con base R}
\porDesarrollar

\semana{4}{Estructura de proyectos, Git y reproducibilidad}
\porDesarrollar

% ------------------------------------------------------------------------------
\subsection*{\color{main}Parte II --- Manipulación}
% ------------------------------------------------------------------------------

\semana{5}{dplyr I: pipe y verbos básicos}
\porDesarrollar

\semana{6}{dplyr II: joins, pivots y manejo de NA}
\porDesarrollar

\semana{7}{Strings, expresiones regulares y \texttt{stringr}}
\porDesarrollar

% ------------------------------------------------------------------------------
\subsection*{\color{main}Parte III --- Programación}
% ------------------------------------------------------------------------------

\semana{8}{Control de flujo y vectorización}
\porDesarrollar

\semana{9}{Funciones: diseño, scoping y debugging}
\porDesarrollar

\semana{10}{Programación funcional: \texttt{purrr} y \texttt{apply}}
\porDesarrollar

% ------------------------------------------------------------------------------
\subsection*{\color{main}Parte IV --- Análisis}
% ------------------------------------------------------------------------------

\semana{11}{Modelado estadístico y series de tiempo}
\porDesarrollar

\semana{12}{Datos a escala: SQL, \texttt{DBI} y \texttt{dbplyr}}
\porDesarrollar

% ------------------------------------------------------------------------------
\subsection*{\color{main}Parte V --- Productos}
% ------------------------------------------------------------------------------

\semana{13}{Visualización con \texttt{ggplot2}}
\porDesarrollar

\semana{14}{Automatización con \texttt{officer} y \texttt{googledrive}}
\porDesarrollar

\semana{15}{Aplicaciones interactivas con Shiny}
\porDesarrollar

\semana{16}{Interfaces de modelos de lenguaje}
\porDesarrollar

% ==============================================================================
\section{Proyecto integrador}
% ==============================================================================

El curso se articula alrededor de un proyecto que los estudiantes construyen a lo largo del semestre, replicando el flujo completo de un proyecto de datos real:

\begin{enumerate}[leftmargin=*]
    \item \textbf{Input:} adquisición de datos desde fuentes diversas.
    \item \textbf{Limpieza:} transformación a formato \textit{tidy} y validación.
    \item \textbf{Modelado:} análisis estadístico apropiado a la pregunta.
    \item \textbf{Producto:} entregable comunicable (documento, dashboard, app).
\end{enumerate}

El proyecto se evalúa por \textbf{checkpoints} alineados con las cinco partes del curso, más una entrega final.

\porDesarrollar{} \emph{(dataset recurrente y entregables específicos por definir)}

% ==============================================================================
\section{Referencias principales}
% ==============================================================================

Las lecturas específicas por sesión se indican en el temario detallado. Las obras de consulta general son:

\nocite{wickham2023r4ds, wickham2019advr, wickham2016ggplot2, wickham2021mastering-shiny}

\printbibliography[heading=none]

\end{document}
